PERSON 3 — PIPELINE, SHAP EXPLAINABILITY, CREDIT SCORING & DECISION SYSTEM
1. Giới thiệu

Sau khi Person 2 hoàn thành quá trình huấn luyện và tối ưu hóa mô hình Machine Learning bằng XGBoost, Person 3 tiếp tục phát triển hệ thống theo hướng triển khai thực tế (deployment-oriented system).

Mục tiêu chính của Person 3 không còn tập trung vào việc cải thiện độ chính xác của mô hình, mà tập trung vào việc xây dựng một hệ thống chấm điểm tín dụng thông minh hoàn chỉnh có khả năng:

Dự đoán rủi ro vỡ nợ
Tạo credit score
Phân loại mức độ rủi ro
Tự động đưa ra quyết định cho vay
Giải thích kết quả dự đoán bằng Explainable AI
Trực quan hóa dữ liệu phục vụ quản trị rủi ro

Hệ thống được xây dựng theo hướng mô phỏng quy trình hoạt động thực tế của ngân hàng và công ty tài chính hiện đại.

2. Kết nối với Person 2

Person 3 sử dụng trực tiếp toàn bộ artifacts đã được huấn luyện từ Person 2.

Các thành phần được load gồm:

Best XGBoost Model
Best Threshold
Model Columns
Pipeline Components

Kết quả:

Best model: XGBoost
Best threshold: 0.71
Number of model columns: 48

Điều này đảm bảo toàn bộ hệ thống của Person 3 kế thừa trực tiếp mô hình tốt nhất đã được tối ưu hóa trước đó.

3. Deployment Pipeline

Person 3 xây dựng deployment pipeline hoàn chỉnh để phục vụ triển khai hệ thống thực tế.

Pipeline được lưu tại:

pipeline/pipeline.pkl

Pipeline bao gồm:

Data preprocessing
Feature alignment
Prediction logic
Credit scoring
Decision engine

Ý nghĩa:

Có thể tái sử dụng trong dashboard/web app
Đảm bảo quy trình prediction đồng nhất
Hỗ trợ triển khai mô hình ngoài môi trường thực nghiệm

Đây là bước cực kỳ quan trọng vì trong thực tế, mô hình Machine Learning cần được đóng gói thành pipeline để triển khai sản phẩm.

4. Dashboard Helper Files

Hệ thống tạo các helper files để hỗ trợ dashboard và deployment:

dashboard/credit_score_logic.py
dashboard/decision_rule.py
dashboard/predict_function.py

Các file này đảm nhiệm:

Tính credit score
Phân loại risk level
Đưa ra quyết định APPROVE / REVIEW / REJECT
Dự đoán khách hàng mới

Điều này giúp hệ thống dễ dàng tích hợp vào ứng dụng thực tế.

5. Batch Credit Scoring

Mô hình được áp dụng lên toàn bộ dữ liệu khách hàng.

Thông tin dataset:

Shape: (32581, 49)
Default rate: 21.82%

Kết quả scoring được lưu:

data/credit_scored_data.csv

Các trường mới được tạo:

probability_of_default
credit_score
risk_level
decision
model_prediction

Ý nghĩa:

Hệ thống giờ đây có thể đánh giá rủi ro cho toàn bộ khách hàng thay vì chỉ dự đoán từng mẫu riêng lẻ.

6. SHAP Explainability

Person 3 triển khai Explainable AI bằng SHAP để giải thích mô hình XGBoost.

Kết quả:

SHAP data type check:
float64    48

Điều này chứng minh toàn bộ dữ liệu đầu vào đã được xử lý đúng định dạng numeric phục vụ explainability.

Các file explainability được tạo:

explainability/shap_summary.png
explainability/shap_bar.png
explainability/shap_waterfall_sample.png
reports/shap_feature_importance.csv
reports/shap_waterfall_sample_data.csv

Ý nghĩa của SHAP:

Giải thích tại sao khách hàng bị đánh giá rủi ro cao
Tăng tính minh bạch cho AI
Hỗ trợ Explainable AI trong lĩnh vực tài chính
Giúp business users hiểu quyết định của mô hình

Các biến ảnh hưởng mạnh nhất tới rủi ro tín dụng gồm:

loan_int_rate
person_income
loan_to_income_ratio
loan_percent_income
person_home_ownership_RENT

Điều này cho thấy:

Lãi suất vay cao
Thu nhập thấp
Tỷ lệ nợ cao

là các nguyên nhân chính làm tăng khả năng vỡ nợ.

7. Visualization & Executive Dashboard

Person 3 xây dựng hệ thống trực quan hóa dữ liệu phục vụ phân tích rủi ro.

Các file được tạo:

visuals/credit_score_distribution.png
visuals/risk_level_distribution.png
visuals/loan_decision_distribution.png
visuals/person3_executive_dashboard.png

Dashboard tổng hợp bao gồm:

Credit Score Distribution
Risk Level Distribution
Loan Decision Distribution
Probability of Default
SHAP Top Features
Business Information Summary

Dashboard giúp:

Theo dõi phân phối rủi ro khách hàng
Hỗ trợ quản trị tín dụng
Giúp nhà quản lý ra quyết định nhanh hơn
Mô phỏng hệ thống Credit Risk thực tế trong ngân hàng
8. Sample Prediction

Hệ thống thực hiện dự đoán thử nghiệm cho một khách hàng.

Kết quả:

{
  "probability_of_default": 0.9966,
  "prediction": 1,
  "credit_score": 302,
  "risk_level": "HIGH",
  "decision": "REJECT"
}

Phân tích:

Xác suất vỡ nợ rất cao: 99.66%
Credit score rất thấp: 302
Hệ thống xếp khách hàng vào nhóm HIGH RISK
Quyết định cuối cùng: REJECT

Điều này cho thấy hệ thống có khả năng hoạt động end-to-end hoàn chỉnh.

9. Ý nghĩa Business

Hệ thống Credit Risk Decision System mang lại nhiều giá trị thực tế.

9.1 Tự động hóa đánh giá tín dụng

Hệ thống có thể đánh giá hàng chục nghìn khách hàng nhanh chóng và nhất quán.

9.2 Explainable AI

SHAP giúp mô hình minh bạch hơn, tăng độ tin cậy khi áp dụng AI trong tài chính.

9.3 Credit Scoring

Credit score giúp business users dễ dàng hiểu kết quả dự đoán thay vì chỉ nhìn probability.

9.4 Decision Support

Hệ thống APPROVE / REVIEW / REJECT hỗ trợ quy trình xét duyệt khoản vay thực tế.

9.5 Khả năng triển khai thực tế

Pipeline và helper files giúp dễ dàng triển khai thành:

Dashboard
Web Application
API Service
FinTech Platform
10. Kết luận

Person 3 đã hoàn thiện giai đoạn triển khai hệ thống Credit Risk theo hướng thực tế.

Các thành phần chính đã hoàn thành:

Deployment Pipeline
Batch Credit Scoring
Decision Engine
SHAP Explainability
Executive Dashboard
Sample Prediction System
Business Visualization

Hệ thống cuối cùng không chỉ là một mô hình Machine Learning đơn thuần mà đã trở thành một nền tảng Credit Risk Analytics hoàn chỉnh.

Dự án kết hợp thành công:

Machine Learning
Explainable AI
Credit Scoring
Risk Analytics
Business Intelligence
Decision Support System

Đây là mô hình rất gần với cách các ngân hàng và công ty FinTech hiện đại triển khai hệ thống đánh giá rủi ro tín dụng trong thực tế.